\section{Introduction}
\label{sec:introduction}

How much critical polygon mass remains after each polygon is weighted by
the square of its length? A bound on the number of large polygons does
not answer this question: a polygon can have many vertices while staying
inside a small region. The useful observable records two visited ports.
Summing their correlation counts a polygon once for each ordered pair
of visits and thus recovers its squared length.

We use the regular honeycomb lattice dual to the triangular lattice of
side length one. A polygon is a simple cycle, with neither orientation
nor distinguished vertex. Its length $|P|$ is its number of vertices,
and its weight is $\rho_{\rm v}^{|P|}$, where
\[
 \rho_{\rm v}=(2+\sqrt2)^{-1/2}.
\]
Translation classes use the translations of the triangular lattice;
the two honeycomb vertex types are consequently retained. Diameter is
Euclidean diameter in this fixed embedding.

\begin{theorem}[Squared-length polygon mass]\label{thm:polygon-moment}
For every $\eta>0$ there is $C_\eta<\infty$ such that, for $H\ge1$,
\[
 \sum_{\substack{[P]\ \operatorname{diam}P\le H}}
       \rho_{\rm v}^{|P|}|P|^2\le C_\eta H^{2/3+\eta}.
\]
The sum is over the simple unoriented polygon translation classes just
defined.
\end{theorem}

Nienhuis predicted the honeycomb critical point and critical exponents
through the dilute $O(n)$ model
\cite{Nienhuis1982}. Duminil-Copin and Smirnov proved the exact
connective constant, hence the value of $\rho_{\rm v}$, using a local
parafermionic identity \cite{DCS2012}.

The spectral parameters used below come from an integrable extension
of this model. Nienhuis obtained soluble square-lattice $O(n)$ models
through the Yang--Baxter equations \cite{Nienhuis1990CriticalMulticritical}.
Ikhlef and Cardy related the local discrete-holomorphicity equations
for parafermionic observables to these integrable weights, with a linear
relation between the spectral parameter and the rhombus angle
\cite[Section~3]{IkhlefCardy2009}. Glazman and Manolescu proved
boundary two-point invariance for columnwise rhombic half-planes
in this Yang--Baxter setting \cite{GM2019}. Zhou and Batchelor studied
finite-size corrections to the bulk transfer-matrix spectra of dilute
integrable models \cite{ZhouBatchelor1997}.
The estimates below also require uniform control of marked positive
lattice sums, including their dependence on the finite geometry and
normalization.

Our main intermediate observable lives on a honeycomb cylinder. Cut the
cylinder along a row of ports in cyclic order
\[
 [p,X,p',Y],\qquad |X|=s,\quad |Y|=b,\quad N=s+b+2\text{ even}.
\]
A port is the midpoint of a triangular edge, hence of a crossed
honeycomb edge. A path between ports includes the initial and final
half-edges and has one factor $\rho_{\rm v}$ for each visited vertex.
The finite transfer uses site logarithms $a_i=2iu_i$, with
$\lambda=\pi/8$ and $\beta=2\lambda=\pi/4$. After a common shift,
the physical logs have heights $0$ or $\beta$. All sites of $Y$ have
height zero; the sites of $X$ can have either height. The two choices
specify the two allowed triangular band directions along the row.
Both marked sites can have either height.

Write $\mathcal C(p,p')$ for the sum of the weights of individual
cylinder polygons visiting both ports. Each polygon is counted once,
including when it winds around the cylinder. The normalized
half-cylinder states defining this sum are described in
Section~\ref{sec:reference}; their physical series converge absolutely.
Let $t$ be the fraction of ordinary $X$ sites at upper height, and put
\[
 \gamma=\frac43,\qquad L_X=\gamma^{-1}\log s,
 \qquad L_Y=\gamma^{-1}\log b.
\]
For fixed $c_1>0$ and a sufficiently small fixed $c_4>0$, the sharp
estimate proved in Theorem~\ref{thm:07_magnetic_analysis-moment} is
\[
 \log\mathcal C(p,p')\le -\frac43\log s
       +\frac{617}{1200}\log(b/s)+o(\log b),
\]
uniformly when $c_1L_Y\le L_X\le L_Y$, $t\in[1/2,1]$, and,
if $t<1$, $s(1-t)\ge e^{c_4L_Y}$. The margin permits an exponentially
small minority of lower sites. Section~\ref{sec:07_magnetic_analysis-trials}
shows that a fixed fraction of the ordered marks on every sufficiently
long planar polygon meet it. This is enough to prove
Theorem~\ref{thm:polygon-moment}; no uniform estimate for every possible
ordered pair is used in that deduction.

There is a second, deliberately coarser, output. Let $A(p,p')$ be the
critical weight of individual simple arcs leaving $p$ on its west
half-edge and arriving at $p'$ on its east half-edge, with the
other marked half-edges vacant and no additional loop components. We prove a fixed
polynomial bound $A(p,p')\le b^C$ for
$L_Y\ge L_X\ge L_*$ when the upper-site fraction of $X$ is at least
$1/2-\delta_1$, for fixed positive $\delta_1,L_*,C$. The ratio
$L_X/L_Y$ may approach zero. This estimate is useful when the cylinder
period grows in one direction while the separation of its marks stays
fixed. Its proof, Theorem~\ref{thm:one-arc}, appears after the sharper two-mark argument and uses
only real integration, so it requires neither the minority-site margin
nor the sharp fluctuation asymptotic.

\subsection{How the proof works}

The finite spin contraction in Section~\ref{sec:06-magnetic-algebra}
selects a single polygon through the two marks. Resolving each ordinary
site against a pair of dual covector families turns that contraction
into a finite sum of contour integrals. The original polygon sum is
positive; the integrals used to estimate it can be complex. Their
normalization is therefore kept exact throughout the argument.

The first normalization issue is independent of the contour saddle.
Section~\ref{sec:denominator} proves a lower bound for the vacuum
polynomial uniform in both site multiplicities. It writes a Pfaffian
ratio as a positive expectation and evaluates its first logarithmic
variation by Cauchy interpolation. This retains the coefficient
$5/24$ even at very unequal multiplicities. The Pfaffian integration
and Cauchy identities are classical tools
\cite{deBruijn1955,IshikawaOkadaTagawaZeng2006}; the uniform estimate
for the present normalization is proved here.

Next, Section~\ref{sec:magnetic-mean} deforms one main contour until its
centering measure is real and positive. The real interaction then
controls a negative Sobolev norm of every species and the $L^2$ norm
of the total primitive. We first verify this on horizontal boundary
lines by an exact finite symbol calculation, then transfer it to the
curved contour by harmonic balayage and a positive Green energy.

Section~\ref{sec:07_magnetic_analysis-trials} uses that coercivity to
bound the particles outside their natural intervals. An imaginary
trial field gains $2L_X$ in the common interval; a different square
completion controls the two shoulders between $L_X$ and $L_Y$.
Their activity sum and the normalization are both retained, rather
than absorbed into a size-dependent constant.

The remaining task is to implement these deterministic trial fields in
the actual integrals. The trial setup introduces the normalized reference
measure in Section~\ref{sec:reference-law}.
Section~\ref{sec:07_magnetic_analysis-fluctuations} computes the
unequal-wall determinant and handles a site mixture tending to one by a
rescaled central coordinate. It then establishes the source and quadratic
bounds for the magnetic reference law and implements the trial fields
by finite-rank particle transport. The sharp source estimate and the
transport adapt the corresponding proofs in the companion on cylinder
weights \cite[Theorem~4.3 and Lemma~5.2]{Electric}; their regularity,
density, and contour hypotheses are verified in the magnetic coordinates.
Finally it combines the trial gain, determinant, reference normalization, explicit
particles and denominator, and sums the resulting positive marked-pair
estimate over planar polygons. The one-arc argument needs less accuracy
and consequently remains valid on its larger parameter domain.
